\documentclass[11pt,a4paper]{article}
\usepackage[margin=2.2cm]{geometry}
\usepackage[T1]{fontenc}
\usepackage{lmodern}
\usepackage{microtype}
\usepackage{xcolor}
\usepackage{hyperref}
\usepackage{booktabs}
\usepackage{longtable}
\usepackage{enumitem}
\definecolor{PilotBlue}{HTML}{155EEF}
\definecolor{PilotInk}{HTML}{15213B}
\hypersetup{colorlinks=true,linkcolor=PilotBlue,urlcolor=PilotBlue}
\setlist{nosep,leftmargin=*}
\title{Pilot Fabric 0.33\\Signed Provider Playback Telemetry}
\author{Tessaris / AION Device Fabric}
\date{30 August 2026}

\begin{document}
\maketitle

\section{Purpose}

Screen perception can establish what appears to be visible, while catalogue
metadata can establish that a provider identifier exists. Neither fact proves
that a household account can play a title, that playback is active, or that a
particular position is current. Pilot Fabric 0.33 introduces a provider-
independent signed telemetry boundary for those stronger claims.

The boundary is designed for authorized native provider applications, owner-
connected account adapters and capable AION device nodes. It does not bypass
closed provider security, scrape credentials or convert a catalogue result into
an entitlement claim.

\section{Trust and scope model}

Each telemetry source is registered locally on the mother brain with:

\begin{itemize}
  \item an Ed25519 public identity;
  \item one provider name;
  \item one private household persona;
  \item a local approval record; and
  \item a least-privilege set of scopes.
\end{itemize}

\begin{longtable}{p{0.27\linewidth}p{0.65\linewidth}}
\toprule
\textbf{Scope} & \textbf{Authority} \\
\midrule
\endhead
\texttt{metadata.read} & Content identifier, title, series, season and episode. \\
\texttt{playback.read} & State, position and duration. \\
\texttt{entitlement.read} & Included, rental, purchase or unavailable status. \\
\bottomrule
\end{longtable}

The scopes are independent. A source with playback access cannot assert an
entitlement, and metadata access cannot assert playback. The shared television
cannot choose a persona or widen a source lease.

\section{Signed envelope protocol}

An adapter emits a canonical JSON
\texttt{pilot.provider-telemetry-envelope.v1} containing source, provider,
persona, nonce, observation time, content, playback and entitlement fields. The
complete envelope is signed with the registered Ed25519 private key.

Fabric rejects a packet when any of the following is true:

\begin{itemize}
  \item the source is unknown or no longer trusted;
  \item the signature is invalid;
  \item the source, provider or persona binding differs;
  \item the observation is more than two minutes old or materially in the future;
  \item the nonce has already been used;
  \item a playback or entitlement state is outside the bounded vocabulary; or
  \item position, duration or field sizes violate their limits.
\end{itemize}

The ingestion endpoint exists only on the trusted LAN HTTPS companion. Its
tokenized path, CSRF check, size limit and rate limiting operate in addition to
the envelope signature; TLS is not treated as a substitute for source identity.

\section{Normalization and privacy}

Pilot stores the normalized fact rather than the provider response. The retained
record contains the bounded content identity, title labels, playback state,
position, duration, entitlement status, evidence hash and source identifier.
Raw responses, account names, email addresses, OAuth credentials, cookies and
provider secrets are excluded.

The source registry is mother-local and owner-readable only. Public adapter keys
are not projected into television, phone or general status snapshots. Audit
records contain telemetry identifiers, hashes, provider, persona, state and the
result of reconciliation, but no credential or raw response.

\section{Fusion with screen perception}

The multi-frame perception timeline accepts signed telemetry as a distinct
source. A signed title can establish provider-verified programme support when
repeated observations agree. A signed provider position outranks an OCR position
and is labelled \texttt{signed\_provider\_adapter}; OCR remains labelled
\texttt{owner\_camera\_ocr} and never becomes provider-verified by implication.

This produces a useful evidence hierarchy:

\begin{enumerate}
  \item signed provider or native-session telemetry;
  \item official catalogue metadata;
  \item permitted television device state;
  \item repeated owner-camera screen observations; and
  \item single-frame OCR or classification evidence.
\end{enumerate}

Conflicting sources remain separate. Pilot does not silently merge a signed
provider title with a different camera observation.

\section{Entertainment execution reconciliation}

A telemetry packet can update only the latest matching execution for the same
persona, provider and content identifier or exact title. A playing state changes
the execution to \texttt{playback\_verified}; a matching non-playing session can
establish only that the provider surface is open. Entitlement changes only when
the source has \texttt{entitlement.read} and supplies a non-unknown state.

Consequently, the following claims remain distinct:

\begin{itemize}
  \item an official route was prepared;
  \item the device received an open command;
  \item the provider surface opened;
  \item the account is entitled to the title; and
  \item the exact title is actively playing.
\end{itemize}

\section{Operational boundary and next work}

Fabric 0.33 completes the secure telemetry contract, verification, persistence,
private-phone display, perception fusion and entertainment reconciliation.
Actual provider coverage depends on supported, owner-authorized adapters. A
service that offers no permitted playback API cannot be treated as if it does.

The next implementation work is to create native television and mobile media-
session adapters where platform policy permits, add authorized provider account
connectors, and field-test disagreement handling across provider telemetry,
webOS state and multi-frame screen perception.

\section{Verification}

The release is covered by tests for signature validation, local source
registration, provider/persona binding, stale and future timestamp rejection,
nonce replay, scope isolation, entitlement boundaries, execution matching,
HTTPS-only ingestion and provider-verified perception timelines. The complete
Fabric regression suite contains 248 passing checks at release time.

\end{document}
